%% ma: L11-B14-0005
%% lop: 11
%% chuong: 4
%% bai: 14
%% dang: 11.14.4
%% loai: TL
%% muc: VD
%% dap_an: "-"
%% nguon: "tai-lieu-hoc-tap-mon-toan-lop-11-theo-chuong-trinh-giao-duc-pho-thong-2018.pdf (Lớp Toán Thầy Duy, Phạm Lê Duy, THPT Chu Văn An), Chương 4 Bài 5, Dạng 3, Ví dụ 5 (tr.524)"
%% kien_thuc: "Lớp 11 Bài 14 (phép chiếu song song)"
%% trang_thai: cho-duyet
%% ly_do: "Dạng mới đề xuất (11.14.x), chờ giáo viên duyệt dạng; tài liệu gốc không có lời giải mẫu, đáp án do kho tự giải và kiểm bằng sympy."
%% ghi_chu: "Đáp số 2 do kho tự giải và kiểm bằng sympy (tài liệu gốc không có lời giải mẫu)."
\begin{ex}
Cho hình hộp $ABCD.A^{\prime}B^{\prime}C^{\prime}D^{\prime}$. Xác định các điểm $M,N$ tương ứng trên các đoạn $AC^{\prime},B^{\prime}D^{\prime}$ sao cho $MN$ song song với $BA^{\prime}$ và tính tỉ số $\dfrac{MA}{MC^{\prime}}$.
\loigiai{\sloppy Đặt hệ vectơ gốc $A$ với $\overrightarrow{AB},\overrightarrow{AD},\overrightarrow{AA^{\prime}}$ là cơ sở: $M=m\overrightarrow{AC^{\prime}}=m(1;1;1)$, $N=B^{\prime}+n\overrightarrow{B^{\prime}D^{\prime}}=(1-n;n;1)$. Khi đó $\overrightarrow{MN}=(1-n-m;n-m;1-m)$ cùng phương $\overrightarrow{BA^{\prime}}=(-1;0;1)$ khi $n=m$ và $\dfrac{1-2m}{-1}=\dfrac{1-m}1$, tức $m=\dfrac23$. Vậy $M$ thuộc $AC^{\prime}$ với $AM=\dfrac23AC^{\prime}$, $N$ thuộc $B^{\prime}D^{\prime}$ với $B^{\prime}N=\dfrac23B^{\prime}D^{\prime}$ và $\dfrac{MA}{MC^{\prime}}=\dfrac{2/3}{1/3}=2$.}
\end{ex}
